% ======================================================================
%  KAVRAM: Agent  (Tier 1)  — kompakt
% ======================================================================

\concept{Orchestration · Orkestrasyon}{Agent}{Hedef verilen modelin kendi adımlarını seçtiği döngü}

\stage{1}{Problem}
Açık uçlu görevlerde her adımı önceden yazamazsın. ``Bunu araştır ve en ucuz
uçuşu bul'' --- atılacak adımlar bulunan şeye göre değişir. Sabit bir akış buna
uyum sağlayamaz; her olasılığı elle kodlamak imkânsızdır.

\keyline{Bazı görevlerin adımları ancak çalışırken belli olur. Bunları önceden senaryolaştıramazsın --- karar verecek bir şeye ihtiyacın var.}

\stage{2}{Neden Var}
Agent, kararı modele bırakır: hedefi al, bir eylem seç (genelde bir araç çağrısı),
sonucu gözlemle, bir sonraki eyleme karar ver, hedef tamamlanana dek döngüye gir.
Açık uçlu görevlerde \emph{esneklik} kazanırsın --- karşılığında \emph{kontrolü}
verirsin.

\stage{3}{Tarihçe}
\begin{timeline}
\yr{2022} ReAct: düşün → eyle → gözlemle döngüsü \tlsep \yr{2023} function calling: yapılandırılmış araç çağrısı \tlsep \yr{2023} otonom agent dalgası \tlsep \yr{2024–26} multi-agent + orkestrasyon; üretimde kontrollü kullanım
\end{timeline}

\stage{4}{Tanım}
\begin{problembox}
\textbf{Agent:} Modelin bir \emph{döngüde} çalıştığı sistem --- hedef verilir;
model bir eyleme (çoğunlukla araç çağrısı) karar verir, sonucu gözlemler ve
hedefe ulaşana dek tekrarlar. Dört bileşen: \textbf{hedef + araçlar + döngü + bellek}.
\end{problembox}

\stage{5}{Zihinsel Model}
\begin{mentalbox}
Workflow bir \textbf{tariftir}: adımlar sabit. Agent ise hedef verilen bir
\textbf{aşçıdır}: ``buzdolabındakiyle akşam yemeği yap.'' Adımları giderken
kendi seçer. Açık uçlu görevlerde daha yeteneklidir --- ama daha az öngörülebilir:
yoldan çıkabilir, döngüde sıkışabilir, bütçeyi yakabilir. Güç ve risk aynı
yerden gelir: kararı ona bıraktın.
\end{mentalbox}

\stage{6}{Mimari Diyagram}
Agent bir döngüdür: düşün → eyle → gözlemle, hedefe dek tekrarla.
\vspace{3pt}
\begin{center}
\begin{tikzpicture}[font=\sffamily\scriptsize, node distance=11mm,
   io/.style={draw=softgray, fill=white, rounded corners=8pt, align=center,
              minimum height=8mm, minimum width=16mm, inner sep=3pt},
   m/.style={draw=accentdk, fill=accent!12, rounded corners=2pt, align=center,
             minimum height=10mm, minimum width=18mm},
   k/.style={draw=accent, fill=wash, rounded corners=2pt, align=center,
             minimum height=10mm, minimum width=18mm},
   flow/.style={-{Stealth[length=2mm]}, accent, line width=0.8pt}]
  \node[io] (goal) {Hedef};
  \node[m, right=12mm of goal] (think) {Düşün\\\textcolor{softgray}{eylem seç}};
  \node[k, right=14mm of think] (act) {Eyle\\\textcolor{softgray}{araç çağır}};
  \node[io, below=10mm of think, fill=good!10, draw=good] (done) {Bitti};
  \draw[flow] (goal) -- (think);
  \draw[flow] (think) -- (act);
  \draw[flow] (act.south) .. controls +(0,-8mm) and +(0,-8mm) ..
        node[below, font=\sffamily\tiny, text=accentdk]{gözlemle} (think.south east);
  \draw[flow] (think) -- node[left, font=\sffamily\tiny, text=good]{hedefe ulaştı} (done);
  \node[font=\sffamily\tiny, text=softgray, above=0mm of act] {döngü};
\end{tikzpicture}
\end{center}
\keyline{Döngü kapanmazsa agent durmaz: bütçe, adım sayısı ve durdurma koşulu \emph{senin} koyduğun sınırlardır.}

\stage{7}{Karar Çerçevesi}
Kitabın en büyük kararı: \textbf{Agent mı, Workflow mu?} Adımlar önceden biliniyor mu?
\vspace{2pt}
\begin{center}
\begin{tikzpicture}[font=\sffamily\scriptsize, node distance=5mm and 8mm,
   q/.style={draw=accent, line width=0.8pt, fill=wash, diamond, aspect=2.4,
             align=center, inner sep=1pt, minimum width=32mm},
   leaf/.style={draw=accentdk, line width=0.8pt, rounded corners=2pt,
                fill=accent!12, align=center, inner sep=4pt, minimum width=22mm},
   wf/.style={draw=softgray, line width=0.8pt, rounded corners=2pt,
                fill=white, align=center, inner sep=4pt, minimum width=22mm},
   yes/.style={good, font=\sffamily\bfseries\scriptsize},
   no/.style={warn, font=\sffamily\bfseries\scriptsize},
   ed/.style={-{Stealth[length=2mm]}, softgray, line width=0.7pt}]
  \node[q] (q1) {Adımlar önceden\\biliniyor mu?};
  \node[wf, below left=13mm and -2mm of q1] (wf1) {\textbf{Workflow}\\öngörülebilir · ucuz};
  \node[q, below right=13mm and -2mm of q1] (q2) {Öngörülebilirlik\\kritik mi?};
  \node[wf, below left=13mm and 4mm of q2, minimum width=20mm] (wf2) {\textbf{Workflow}\\(+ dar agent)};
  \node[leaf, below right=13mm and 4mm of q2, minimum width=20mm] (ag) {\textbf{Agent}\\+ sınır + eval};
  \draw[ed] (q1) -- node[yes,left]{Evet} (wf1);
  \draw[ed] (q1) -- node[no,right]{Hayır} (q2);
  \draw[ed] (q2) -- node[yes,left]{Evet} (wf2);
  \draw[ed] (q2) -- node[no,right]{Hayır} (ag);
\end{tikzpicture}
\end{center}
\keyline{Varsayılan Workflow olsun. Agent'ı yalnızca adımlar gerçekten önceden bilinemiyorken seç.}

\stage{8}{Üretim Senaryosu}
\begin{prodbox}[Üretim Senaryosu · Agent'tan Workflow'a Geri Çekilmek]
\textbf{Bağlam.} Bir ekip her şeyi tek bir otonom agent'la çözmeye çalışıyor.
Maliyetler öngörülemez, hata ayıklama imkânsız, bazı koşular dakikalarca döngüde
kalıyor.

\smallskip
\textbf{Çözüm.} Vakaların \%90'ının aslında \emph{sabit} adımlardan oluştuğu
görülür. Bunlar deterministik bir workflow'a alınır; agent yalnızca gerçekten
açık uçlu \%10'luk ``kuyruk'' için, sıkı bütçe ve adım sınırıyla saklanır.

\smallskip
\textbf{Sonuç.} Maliyet öngörülebilir, hata ayıklama kolay, kalite artar ---
agent'ın esnekliği yalnızca gerçekten gerektiği yerde ödenir.
\end{prodbox}

\stage{9}{Anti-pattern}
\begin{antibox}[Anti-pattern · ``Her Problem Bir Agent'']
Deterministik bir workflow'un daha ucuz, hızlı ve güvenilir olacağı yerde otonom
agent kullanmak. Sabit, basit görevler için agent: öngörülemez maliyet, zor hata
ayıklama, döngüde sıkışma.

\smallskip
\textbf{Belirti:} ``Basit bir iş için agent kurduk; maliyet ve davranış kontrolden çıktı.''\\
\textbf{Kök neden:} Esnekliği her yerde istemek; kontrolün değerini görmemek.\\
\textbf{Düzeltme:} Varsayılanı workflow yap; agent'ı yalnızca açık uçlu kuyruğa, sınırlarla uygula.
\end{antibox}

\stage{10}{Alternatifler \& Ödünleşimler}
\vspace{2pt}
\begin{tabularx}{\linewidth}{@{}>{\raggedright\arraybackslash}p{26mm} L L@{}}
\arrayrulecolor{hairline}\toprule
{\sffamily\bfseries\small Yaklaşım} & {\sffamily\bfseries\small Güçlü} & {\sffamily\bfseries\small Ödünleşim} \\
\midrule
\textbf{Agent} & Açık uçlu, uyarlanır görev & Öngörülemez; pahalı; zor debug \\
\addlinespace[2pt]
\textbf{Workflow} & Öngörülebilir, ucuz, denetlenir & Sabit; uyum sağlamaz \\
\addlinespace[2pt]
\textbf{Tek prompt} & En basit görev & Çok adımlı/araçlı işe yetmez \\
\addlinespace[2pt]
\textbf{Prompt chain} & Bilinen birkaç adım & Dallanma/karar gerekirse yetersiz \\
\bottomrule
\end{tabularx}

\stage{11}{İlişkili Kavramlar}
\term{Workflow} (karşıtı; varsayılan) \tlsep \term{Tool Calling} (eylem mekanizması) \tlsep
\term{Reasoning} (planning/think adımı) \tlsep \term{Memory} (uzun görev belleği) \tlsep
\term{Multi-Agent} (birden çok agent) \tlsep \term{Router} (basit dağıtım)

\stage{12}{Özet}
\begin{itemize}
  \item Agent: hedef + araçlar + döngü + bellek; model kendi adımlarını seçer.
  \item Esneklik kazandırır, kontrolü alır --- güç ve risk aynı yerden gelir.
  \item Adımlar önceden biliniyorsa Workflow; ancak gerçekten açık uçluysa Agent.
  \item Sınır senin sorumluluğun: bütçe, adım sayısı, durdurma koşulu, eval.
  \item ``Her problem bir agent'' kitabın en pahalı anti-pattern'idir.
\end{itemize}

\stage{13}{Mühendislik İlkesi}
\begin{principlebox}
\textbf{İlke \#11 — Kontrol edebildiğin her şeyi workflow yap; agent'ı yalnızca mecbur kalınca kullan.}\\[4pt]
\normalsize Otonomi bir maliyettir, bir özellik değil. Yalnızca adımları önceden
bilemediğin gerçek belirsizlik için öde.
\end{principlebox}
